\chapter{Khởi Động Cho Giờ Chủ Nhiệm}
\chapterintro{Mở giờ sinh hoạt sao cho lớp chịu nói thật}{Giờ chủ nhiệm cần một đầu vào mềm, tích cực và an toàn để cả lớp không bật chế độ phòng thủ ngay.}

\section{Một điều lớp tự hào trước đã}
\begin{khoidongbox}{Hoạt động khởi động}
Mở bằng câu hỏi: \textit{``Tuần này lớp mình tự hào điều gì trước?''}
\end{khoidongbox}
\begin{visaohieuqua}
Khen trước không phải để né lỗi, mà để lớp bước vào buổi sinh hoạt với tinh thần hợp tác hơn.
\end{visaohieuqua}
\begin{cachlam5phut}
Lấy 2 đến 3 điểm sáng rồi mới sang phần cần sửa. Không mở giờ chủ nhiệm bằng bảng lỗi ngay lập tức.
\end{cachlam5phut}
\ngancach

\section{Một điều lớp muốn bớt}
\begin{khoidongbox}{Hoạt động khởi động}
Hỏi thẳng nhưng nhẹ: \textit{``Tuần tới lớp mình muốn bớt điều gì nhất?''}
\end{khoidongbox}
\begin{visaohieuqua}
Câu hỏi gom sự chú ý vào một việc có thể thay đổi thật. Nó giúp giờ sinh hoạt bớt lan man và bớt lên án.
\end{visaohieuqua}
\begin{cachlam5phut}
Chỉ chốt một điều, viết lên bảng, rồi quay lại nó ở cuối buổi sinh hoạt.
\end{cachlam5phut}
\ngancach

\section{Chọn mood cho tuần mới}
\begin{khoidongbox}{Hoạt động khởi động}
Cho lớp chọn giữa vài mood đơn giản cho tuần mới: bình tĩnh hơn, đúng giờ hơn, vui hơn, tập trung hơn.
\end{khoidongbox}
\begin{visaohieuqua}
Học sinh dễ nhập cuộc khi được chọn một trạng thái mong muốn thay vì chỉ nghe nhắc lỗi cũ.
\end{visaohieuqua}
\begin{cachlam5phut}
Cho giơ tay chọn mood, rồi hỏi 2 em vì sao muốn lớp đi theo hướng đó.
\end{cachlam5phut}
\ngancach

\section{Lời hứa một dòng}
\begin{khoidongbox}{Hoạt động khởi động}
Mời mỗi tổ chốt một lời hứa trong đúng một dòng cho tuần tới.
\end{khoidongbox}
\begin{visaohieuqua}
Một dòng thì dễ nhớ, dễ nhắc, dễ làm hơn những cam kết dài và đẹp nhưng mờ.
\end{visaohieuqua}
\begin{cachlam5phut}
Mỗi tổ chỉ một câu. Chọn ra một lời hứa chung của cả lớp để giữ trong tuần.
\end{cachlam5phut}
\ngancach

\section{Một câu cảm ơn trong lớp}
\begin{khoidongbox}{Hoạt động khởi động}
Cho học sinh nói một câu cảm ơn ngắn dành cho bạn, tổ, hoặc cả lớp.
\end{khoidongbox}
\begin{visaohieuqua}
Lời cảm ơn làm giờ chủ nhiệm mở mềm hơn rất nhiều. Nó cũng kéo lớp về tinh thần xây tập thể thay vì chỉ xử lý sự vụ.
\end{visaohieuqua}
\begin{cachlam5phut}
Lấy 3 đến 5 câu là đủ. Sau đó mới chuyển sang phần phản ánh hoặc điều chỉnh.
\end{cachlam5phut}